\section{Numerical method}
\label{sec:method}

\status{Ready to write: every part described here is implemented and verified.
The subsection on adaptive refinement (\ref{sec:amr}) is the exception --- the
criterion exists and is measured in two dimensions, but has not been ported to
three.}

\subsection{Shared machinery, and the single point of difference}
\label{sec:shared}

This subsection carries the methodological claim of the paper, so it states the
arrangement precisely.

Both methods use the same projection method, the same variable-coefficient
momentum solve, and the same property interpolation. The front-tracking
coupling is a function pointer, null by default, invoked inside the multiphase
step \emph{after} the source terms and \emph{before} the viscosity and density
interpolation. At that point it overwrites the volume fraction from the
geometry of the front, so the properties are built from the new fraction. VOF
reaches the same field at the same point by advecting (\ref{eq:transport}) with
a geometric PLIC reconstruction.

Everything downstream of that point --- $\rho(\mathbf{x})$, $\mu(\mathbf{x})$,
the momentum solve, the pressure projection, the output probes --- is
bit-identical code. The consequence is that a measured difference between the
two is a difference in how the interface is represented, and not in the solver,
the mesh, the time step, the initial geometry, or the instrument.

\subsection{VOF: geometric PLIC}

\missing{PLIC reconstruction and the split advection --- short, since it is the
pre-existing method; cite the reference implementation}

\paragraph{Initial volume fraction.}
The initial fraction is computed by \textbf{sub-cell subdivision}: each
interface cell is split into $16^3$ sub-cells, each resolved by the closed-form
plane--box intersection volume, and the sum is condensed onto the coarse cell.
This matters more than it appears. A one-cell smoothed step --- the obvious
shortcut --- gives a bubble volume error of \num{1.23e-2} at the running
resolution against \num{6.11e-5} for the subdivision, and its measured order of
convergence is erratic (1.59, 2.90, 0.99) because the ramp width is tied to $h$.
Since VOF conserves faithfully whatever mass it is given, an initial-condition
error of that size never recovers.

\subsection{Front-tracking: explicit triangulated surface}

The interface is a closed triangulated surface with explicit connectivity,
built as an icosphere by recursive subdivision of an icosahedron.

\paragraph{Surface tension by line integral.}
The force on each triangle is obtained from the line integral along its edges,
$\mathbf{f} = \sigma(\mathbf{b}-\mathbf{a})\times\mathbf{n}$, distributed to the
two endpoints. The force is $-\sigma\nabla A$ --- tension minimises area ---
and the sign was derived rather than guessed, then checked against the numerical
area gradient.

\paragraph{The Laplace oracle in three dimensions.}
For a sphere the pressure jump is $\Delta p = 2\sigma/R$, \emph{not} $\sigma/R$.
The two-dimensional value is $\sigma/R$, and a 2D implementation that closes to
0.02\% against it would, if the number were carried over, validate the wrong
formula. The distinction is flagged in the header, the harness and the example.

\paragraph{Surgery.}
Split, collapse and flip, with a curved midpoint
\begin{equation}
\mathbf{m} = \frac{\mathbf{a}+\mathbf{b}}{2}
  - \frac{1}{8}\Big[\big((\mathbf{b}-\mathbf{a})\cdot\mathbf{n}_a\big)\mathbf{n}_a
  + \big((\mathbf{a}-\mathbf{b})\cdot\mathbf{n}_b\big)\mathbf{n}_b\Big],
\end{equation}
which is the PN-triangle (Hermite) construction. Surgery keeps the triangle
edge matched to the cell size as the bubble deforms.

\paragraph{Volume fraction from the surface.}
The fraction is obtained from the signed distance and normal to the surface,
then from the closed-form plane--box volume. The sign comes from parity
ray-casting, never from the nearest triangle's normal --- the latter is wrong
for a cell straddling a concavity.

\subsection{Interface extraction for comparison}
\label{sec:mc}

VOF has no interface as an object; the comparison needs one. The isosurface
$f = 1/2$ is extracted by marching cubes on the \emph{dual} grid --- the volume
fraction lives at cell centres, so the lattice nodes are cell centres --- with
vertices \textbf{welded} by lattice edge, which makes the extracted surface
continuous. The alternative, writing the per-cell PLIC planes, gives a surface
that is discontinuous from one cell to the next; that discontinuity is the right
object for auditing the reconstruction and the wrong one for measuring shape.

Both methods therefore emit a triangulated surface in the same file format, and
every shape quantity in this paper is computed from those surfaces by the
\emph{same} routine. The instrument is common; only the surface differs.

\paragraph{On adaptive meshes the extraction runs on the finest level only.}
This is sufficient, and not by economy: the band criterion refines to the finest
level with a floor of five fine cells on each side of the interface, and outside
that band the fraction is exactly 0 or 1. The eight-node stencil therefore never
reaches the seam between levels, which is what normally breaks isosurface
extraction on adaptive meshes.

\subsection{Adaptive refinement}
\label{sec:amr}

\status{The criterion is implemented and measured in two dimensions; it has not
been ported to three. The 2D signature is planar and the seed spatial index
needs the third extent.}

\missing{the criterion: maximum of a geometric band around the interface and a
vorticity threshold, with hysteresis --- a cell already refined only coarsens
once $|\omega|$ falls below half the threshold. Port to 3D pending.}
